\section{Interface Surface and the Governed-Asset Contract}
\label{app:api}

\subsection{The shape every registry presents}

\statRouteDecls{} route declarations across \statRouteModules{} modules is too many
to enumerate, and enumerating them would not tell a reader anything structural. What
is structural is that each family's registry presents the same shape, and that the
\emph{semantics} behind identical-looking operations differ per family exactly as
\tabref{tab:families} records.

\begin{lstlisting}[caption={The shape shared by every family's registry surface.
Identical operation names; per-family semantics. The comments mark where
\tabref{tab:families} is the authority.},label={lst:contract},
captionpos=b,language={}]
GET    /{family}                 list, project-scoped
POST   /{family}                 create a typed definition
GET    /{family}/{id}            read the current record
PATCH  /{family}/{id}            mutate  -- scope-checked, audited
GET    /{family}/{id}/versions   history -- SEMANTICS VARY (Tab. 1)
GET    /{family}/{id}/diff       compare two versions

# Present only where the family implements the facet:
POST   /{family}/{id}/test       bounded run against fixtures
POST   /{family}/{id}/evaluate   scorecard over assembled evidence
POST   /{family}/{id}/calibrate  enqueue a refinement job
POST   /{family}/{id}/promote    rotate the live target -- operator/admin
POST   /{family}/{id}/rollback   restore the recorded prior target
\end{lstlisting}

The lower block is the interesting one. Its endpoints exist per family only where
that family wires the corresponding facet, which makes the interface surface itself
a reflection of \secref{sec:arch:claim}: availability is structural, enforcement is
wired. A client cannot discover a promote endpoint on a test set, because there is
none to discover.

\subsection{What a promote must record}

The prompt-alias service records the following fields before its external effect
(\algref{alg:promote}). Authoring routes cannot rotate an alias; the direct
promotion, rollback, assistant-publish, and refinement \mode{Apply} call sites use
the shared service. This is a verified prompt-family contract, not a system-wide
invariant for every resource family (\secref{sec:limitations}).

\begin{enumerate}[leftmargin=1.5em]
\item \textbf{The outgoing target.} What the \livetarget{} pointed to before the
      rotation. Without this, rollback is a guess.
\item \textbf{The incoming version identity.} Which version is becoming live, by
      immutable identifier rather than by label.
\item \textbf{The evidence reference.} A refinement Apply records its approval,
      candidate/evaluation snapshots, and gate source; a direct operator release
      records the advisory gate state and any attributed override. Missing evidence
      is represented as \code{none}, not silently omitted.
\item \textbf{The authorizing principal and scope.} Who acted, under which of the
      \statScopesW{} scopes.
\item \textbf{An audit row}, written on the caller's session where the mutation is
      database-resident (\secref{sec:components}).
\item \textbf{A release operation state}, progressing through
      \code{prepared}, \code{applying}, and \code{applied}, or remaining visible
      as \code{failed}/\code{reconcile\_required}.
\end{enumerate}

\subsection{Admission, in the order it happens}

Order matters here, and stating it prevents a specific misreading of the body cap
described in \secref{sec:components}.

\begin{enumerate}[leftmargin=1.5em]
\item For a protected published-service invoke, the Bearer token is validated
      \emph{before} the request body is consumed. An unauthenticated caller cannot
      make the server read a large body.
\item For both public and protected invokes, raw request chunks are then counted
      against the configured maximum during bounded parsing.
\item Policy and token are revalidated under the enqueue lock, so a policy change
      between admission and enqueue cannot be raced.
\end{enumerate}

This is a per-request application bound. It is not connection-level or IP-level
flood protection, and step~1 does not change that --- it only ensures the bound is
applied to authenticated callers only where a token is required.
